\documentclass[a4paper,10pt,preprint]{elsarticle}

% --- UNIVERSAL PREAMBLE BLOCK ---
\usepackage[a4paper, top=2.5cm, bottom=2.5cm, left=2cm, right=2cm]{geometry}

\usepackage[english, bidi=basic, provide=*]{babel}

\babelprovide[import, onchar=ids fonts]{english}

% Set default/Latin font to Sans Serif in the main (rm) slot

% Packages required for the table and general formatting
\usepackage{float}      % For the [H] placement modifier
\usepackage{tabularx}   % For automatically wrapping columns to a target width
\usepackage{lineno,hyperref}
\usepackage{booktabs}
\usepackage{listings}
\usepackage{xcolor}
\usepackage{graphicx}
\usepackage{longtable}
\usepackage{amsmath}

\begin{document}

% We place this before the frontmatter to guarantee it appears on the very first page 
% as explicitly requested, before the title and authors.
  \begin{figure}[H]

\centering

\includegraphics[width=1.0\linewidth]{figures/graphical_abstract.png}

\caption{End-to-End Multi-Modal SciML Dosimetry Pipeline. Phase 1 shows raw data of different modalities in coronal view with typical acquisition misalignments (rotation and scaling discrepancies). Phase 2 demonstrates the result of unified preprocessing and batch alignment on a shared grid. Phase 3 illustrates the integration with Julia framework. Phase 4 displays possible use case scenario with UDE based  dosimetry prediction}

\label{fig:flowchart}

\end{figure}

\begin{frontmatter}

% Title
\title{MedImages.jl: A High-Performance, Differentiable Julia Framework for Medical Image Processing and Metadata Preservation}

% Properly using elsarticle author, address, and corresponding author macros
\author[1]{Jakub Mitura\corref{cor1}}
\ead{jakub.mitura@med.ovgu.de}

\author[2]{Divyansh Goyal}
\author[1]{Michael C. Kreissl}
\author[1]{Joanna Wybranska}

% Affiliations
\address[1]{Division of Nuclear Medicine, Department of Radiology \& Nuclear Medicine, Faculty of Medicine, Otto von Guericke University Magdeburg, 39120 Magdeburg, Germany}
\address[2]{Department of Artificial Intelligence and Machine Learning,
Guru Gobind Singh Indraprastha University, Delhi, India}

% Corresponding Author Text
\cortext[cor1]{Corresponding author. Full Postal Address: Leipziger Str. 44, 39120 Magdeburg, Germany. Telephone: [Insert Telephone Number Here]}

\begin{abstract}
Rationale: Imaging data driven artificial intelligence  (AI) research in Nuclear Medicine (NM) increasingly depends on large-scale multi-modal 3D/4D datasets that must be resampled, aligned and quantitatively normalized for training of AI models without losing the spatial and clinical metadata. Existing pipelines often struggle with (1) biobank-scale volume with  I/O overhead and execution speed, (2) metadata drift when images are converted to raw arrays, and (3) end-to-end differentiability for novel physics-informed machine learning.

\textbf{Methods:} We developed \texttt{MedImages.jl}, a native Julia library that binds voxel data to spatial metadata (origin, spacing, direction) and clinical/temporal fields using a typed \texttt{MedImage} structure and a batched multimodal container (\texttt{BatchedMedImage}) for synchronized transforms. Volume-scale performance was evaluated on a 100-subject PET/CT alignment pipeline (target grid $512^3$) comparing HDF5-based persistence with a MONAI Python caching pipeline. Execution speed for resampling, orientation changes, and fused affine transforms was benchmarked on \(256 \times 256 \times 128\) volumes using MedImages.jl GPU, MedImages.jl CPU, and SimpleITK CPU implementations. Metadata integrity was assessed by SUVbw agreement with SimpleITK and by an end-to-end SPECT/CT transformation chain with liver-mask-based SUV and volume readouts. Differentiability was validated via gradient tests and learned spatial pipelines using inverse-rotation spatial transformer and further extended to a physics-informed neural-network (PINN) dosimetry model for \(^{177}\)Lu-PSMA theranostic workflows, trained against semi-Monte Carlo (sMC) reference dose maps (DM) and compared against analytical baselines and state-of-the-art neural-network (NN) for voxel-level dosimetry.

\textbf{Results:} \texttt{MedImages.jl} achieved a \(7.2\times\) turnaround speedup over the MONAI caching pipeline in the 100-subject benchmark and, for isolated spatial operations, MedImages.jl GPU outperformed MedImages.jl CPU by \(115\times\), \(71\times\), and \(135\times\) for resampling, orientation changes, and fused affine transforms, in that order. The MedImages.jl CPU additionally outperformed the SimpleITK CPU by \(17\times\), \(31\times\) and \(8.1\times\) in spatial preprocessing operations, respectively. SUV\textsubscript{bw} matched SimpleITK to approximately \(10^{-14}\) precision, and the end-to-end transformation chain preserved liver-mask-based quantitative readouts across the 10-patient cohort, with a mean SUV deviation of \(0.17\%\) (range: \(0.00\%\)--\(0.40\%\)) and a liver volume deviation of \(0.10\%\) (range: \(0.01\%\)--\(0.28\%\)), consistent with interpolation effects rather than metadata deterioration. Learned spatial transforms reduced reconstruction error by \(>65\%\), and UDE dosimetry refinement converged to an MSE of \(4.14\times10^{-4}\) versus sMC ground-truth dose maps (DM), achieving the best overall performance among the evaluated NN.

\textbf{Conclusion:} \texttt{MedImages.jl} provides a unified, GPU-accelerated, metadata-aware, and differentiable foundation that addresses scalability, speed, differentiability, and metadata integrity for quantitative nuclear medicine preprocessing and AI/SciML workflows.

\end{abstract}

\end{frontmatter}

\textbf{Key words:} nuclear medicine; PET/CT; SUV; metadata; differentiable programming

\section{Introduction}

Nuclear medicine diagnostic data, such as PET and SPECT, provide quantitative measurements of biologic processes that support oncologic staging, assessment of therapy response, and planning of theranostic treatment. Correct processing and interpretation of functional imaging data require that measured uptake is assigned to the same patient-space coordinates as the anatomic reference (CT or MRI). Hybrid imaging therefore combines complementary modalities within a single, physically defined sampling space specified by image origin, grid or voxel spacing, and orientation. 

% NM diagnostic data, such as PET and SPECT, provide quantitative measurements of biologic processes that support oncologic staging, assessment of therapy response, and theragnostic treatment planning. Correct processing and interpretation of functional imaging data require that measured uptake is assigned to the same patient-space coordinates as the anatomic reference (CT or MRI). Hybrid imaging therefore combines complementary modalities within a single, physically defined sampling space specified by image origin, grid or voxel spacing, and orientation. %

As multimodal imaging is increasingly reused as input for AI models, managing data volume and maintaining heterogeneous acquisition metadata has become a central practical challenge \cite{Simon_2024}. A high computational data preparation process typically requires: loading of a large number of examinations from disk, mapping each modality into a shared patient-space coordinate system, resampling onto standardized grids and applying geometric transforms for augmentation or registration. Preserving spatial and quantitative metadata is necessary throughout the multistep preprocessing phase so that voxel values remain physically interpretable. 

%As multimodal imaging is increasingly reused as input for AI models, managing data volume and maintaining heterogeneous acquisition metadata has become a central practical challenge \cite{Simon_2024}. High computational data preparation processes typically require: loading of a large number of examinations from disk, mapping each modality into a shared patient-space coordinate system, resampling onto standardized grids and applying geometric transforms for augmentation or registration. Preserving spatial and quantitative metadata is necessary throughout the multistep preprocessing phase so that voxel values remain physically interpretable. %

Metadata management is not only a computational but also a safety-relevant issue. Converting medical image objects (DICOM, NIfTI) to raw arrays for deep learning can decouple voxel intensities from their patient-space definition and from clinical metadata used for quantification. This ``metadata drift'' can yield visually plausible results that are spatially inconsistent or quantitatively invalid \cite{Elias_2023}.

In widely used image-processing environments such as SimpleITK \cite{LowekampChen2013} and MONAI \cite{CardosoLi2022}, data handling is often built around workflows and chains of separate libraries rather than architectures tailored to biobank-scale diagnostic imaging. As a result, a substantial fraction of total processing time is used not only for core preprocessing operations itself, but also to repeatedly copy very large image volumes between intermediate storage steps and different software components, and to ensure the consistency of spatial information \cite{Dexl_2025}.

These limitations become even more pronounced in emerging Scientific Machine Learning (SciML) workflows, which increasingly rely on end-to-end optimization in which geometric preprocessing participates in training of AI models. In NM, such workflows include learned registration and motion correction, data augmentation, and physics-informed models for image reconstruction or dosimetry. Although Python-based frameworks provide differentiable operations, practical composability across heterogeneous software components often requires committing the workflow to a single tensor ecosystem, creating friction when integrating third-party tools or classical numerical solvers \cite{PalBhattacharya2024, EschleGl2023}. 

To address these limitations, we developed MedImages.jl, a Julia-based imaging ecosystem that converts high-level code into efficient machine instructions using just-in-time intermediate optimization and enables native implementations of geometric image operations. As a result, these operations can run efficiently on both CPUs and GPUs  \cite{BezansonChen2018}. 

MedImages.jl binds voxel data to spatial and acquisition metadata through typed image representations, so that the typical for preprocessing geometric operations update the underlying physical description rather than operating on detached arrays. Core operators like resampling, affine warping, orientation handling, cropping and padding are implemented as native Julia kernels with transparent GPU acceleration, reducing practical bottlenecks in repeated 3D processing and supporting multimodal synchronization through batched representations.

The presented framework is designed for interoperability with the broader Julia SciML ecosystem and with automatic differentiation libraries such as Zygote.jl and Enzyme.jl, which allows for differentiable spatial pipelines \cite{10.1145/3458817.3476165}. This permits geometric image operations to be incorporated directly into end-to-end optimization, allowing spatial transformations to influence model training rather than remain external preprocessing step. Julia further reduces software fragmentation by enabling image transformations, reconstruction algorithms, numerical solvers, and optimization routines to operate within a single coherent framework instead of across separate software layers that must be manually linked \cite{RoeschGreener2023}. Since these compiler-level differentiation engines can act on arbitrary Julia code, gradients can propagate through solvers and GPU kernels without rewriting the underlying workflow in a specialized tensor framework \cite{Fischer\_2019}. 

In this work, we evaluate MedImages.jl across NM–relevant benchmarks and experiments, both in isolation and in integration with complementary Julia libraries, with emphasis on (i) scalability and execution speed  for large PET/CT cohorts with high preprocessing demands, (ii)  metadata preservation to maintain the physical interpretability required for quantitative NM, and (iii) differentiability of geometric operations for learned spatial transforms.


\section{Methods}

MedImages.jl was developed to address several practical requirements that are highly relevant in medical imaging, including preservation of image geometry and acquisition metadata, reliable handling of clinical image formats, efficient processing of large multimodal datasets, and support for optimization-based learning workflows. 
A central design feature is the \verb|MedImage| data structure, which stores voxel values together with essential spatial information such as image position, voxel spacing, orientation, patient identity, and acquisition time. This avoids a common problem in AI-oriented workflows, where images are converted into generic tensors and thereby lose information needed to interpret them in physical space \cite{}. 
For robust reading and writing of clinical formats such as DICOM and NIfTI, MedImages.jl uses Insight Segmentation and Registration Toolkit (ITK) through a wrapper interface, whereas spatial operations such as rotation, translation, and resampling are implemented directly in Julia language. This design combines the reliability of a mature medical imaging toolkit for file handling with the flexibility to use spatial operations in AI learning pipelines. 
For larger studies, the \verb|BatchedMedImage| structure allows multiple images to be processed together while preserving the spatial properties of each individual volume. In combination with HDF5-based storage and GPU acceleration, this architecture improves scalability, and supports workflows in which image transformations can be incorporated directly into AI model optimization.  
Moreover, compatibility with the broader Julia SciML ecosystem allows MedImages.jl to interoperate with other scientific computing libraries.

An overview of  the core framework functionalities is summarized in Table \ref{tab:medimages_functionalities}.

To demonstrate that MedImages.jl addresses the central computational and metadata-related challenges of large-scale medical image analysis, we conducted three proof-of-concept experimental series. Additional implementation details, dataset mappings, model architectures, baseline definitions, and example code are provided in the Appendix. 

\begin{table}[H]
\centering
\caption{Core functionalities and data structures of MedImages.jl for 3D medical image processing.}
\label{tab:medimages_functionalities}
\renewcommand{\arraystretch}{1.2}
\begin{tabular}{|p{0.24\textwidth}|p{0.68\textwidth}|}
\hline
\textbf{Functionality} & \textbf{Description} \\
\hline
\texttt{MedImage} data structure & Unified medical image representation that stores voxel data together with clinical and vectorized spatial metadata. \\
\hline
\texttt{BatchedMedImage} structure & Batched representation for joint processing of multiple images while preserving the spatial properties of each individual volume. \\
\hline
Imaging-format-compatible I/O & Loading of standard medical image formats, including NIfTI and DICOM, via \texttt{ITKIOWrapper}. \\
\hline
HDF5 support & Fast loading and saving in big-data compatible format, facilitating efficient intermediate storage in repeated processing and deep-learning workflows. \\
\hline
GPU acceleration & Transparent GPU support for spatial transformations through \texttt{KernelAbstractions.jl}, enabling accelerated parallel execution of preprocessing and geometric operations. \\
\hline
Resampling & Native Julia functions for resampling, enabling voxelwise correspondence across modalities. \\
\hline
Spatial transformations & Native Julia functions for manipulation of spatial properties of 3D image volumes, including \texttt{rotate\_mi}, \texttt{translate\_mi}, \texttt{scale\_mi}, \texttt{resample\_to\_spacing}, \texttt{crop\_mi}, and \texttt{pad\_mi}. \\
\hline
Orientation management & Tools for handling, preserving, and converting image orientations across neuroimaging and medical imaging coordinate systems. \\
\hline
Automatic differentiation & Differentiable implementation of spatial transformations with well-defined gradients through \texttt{Zygote.jl} and \texttt{Enzyme.jl}. \\
\hline
\end{tabular}
\end{table}




\textbf{Experimental Series 1: Evaluation of Scalability and Execution Speed}

For large-scale imaging studies, we evaluated and compared MedImages.jl with contemporary frameworks (MONAI and SimpleITK) with respect to two key requirements: scalability to multimodal cohorts and runtime of core spatial preprocessing operations in repeated deep-learning workflows.

To assess scalability and computational overhead, we designed a 100-case preprocessing and training benchmark based on the autoPET dataset \cite{Dexl_2025}. This setup was intended to reflect a realistic large-cohort scenario in which multiple PET/CT studies must be repeatedly loaded and harmonized during iterative model training.
MedImages.jl was compared with MONAI using its PersistentDataset mechanism \cite{KaliyugarasanLundervold2023}, thus assessing not only the computational performance at the framework-level, but also the effect of different metadata persistence strategies. 
The main preprocessing task was multimodal spatial alignment of heterogeneous PET and CT volumes on a common isotropic reference grid of 512 × 512 × 512 voxels. This harmonization is essential for multimodal neural networks, which require strict voxelwise spatial correspondence across input channels. In practical terms, alignment means transforming both modalities into the same physical coordinate system such that a voxel in position (i,j,k) of the PET volume corresponds to the same anatomical location as the voxel in position (i,j,k) of the CT volume.
Data handling performance was evaluated under repeated training conditions, beginning from epoch 2 onward, under the assumption that the data had already undergone initial preprocessing steps and had been stored in cache.

To assess raw computational speed independently of the caching benchmark, we additionally measured isolated spatial operations across multiple scales. These tests included resampling, orientation changes, and affine transformations, which represent some of the most common geometric operations in multimodal medical image preprocessing. Benchmarks were performed on single volume of size 256 × 256 × 128.  SimpleITK was evaluated on CPU, reflecting its typical use for such geometric operations, whereas MedImages.jl was evaluated on both CPU and GPU in order to quantify the effect of hardware acceleration on the same class of tasks.

Both the 100-subject caching pipeline benchmark and the isolated raw computational speed benchmarks (resampling, orientation changes, fused affine transformations) were executed on same hardware configurations: an Intel Core Ultra 9 285K processor for CPU-bound computations and an NVIDIA RTX 3090 GPU for hardware-accelerated processing."


\textbf{Experimental Series 2: Verification of Metadata Integrity and Quantitative Clinical Fidelity  }

We verified the quantitative accuracy of MedImages.jl by comparing its Standardized Uptake Value (SUV) calculations against the gold standard SimpleITK across the 10-patient cohort. 
End-to-end consistency experiment was explicitly performed on a cohort of 10 clinical subjects to verify that metadata preservation and quantitative fidelity are maintained throughout complex spatial transformation pipelines. For each patient, automated liver masks were generated using \textit{TotalSegmentator} \cite{Wasserthal2023} as spatial references. The evaluation involved a multi-modal sequence consisting of (1) converting SPECT intensities to SUV, (2) resampling the functional data to the native CT target space to ensure baseline alignment, (3) changing the spacing to a 2.0\,mm isotropic grid, (4) applying a 45-degree rotation around the Z-axis, and (5) applying a 10-voxel translation along the X-axis (see Figure \ref{fig:transformations}). For the segmentation masks, nearest-neighbor interpolation was strictly applied to prevent boundary blurring, while linear interpolation was utilized for the continuous intensities.
Quantitative evaluation was performed by measuring mean SUV and total lesion volume  (TLV) within the segmentation mask before and after the transformation sequence.

\begin{figure}[H]
\centering
\includegraphics[width=1.0\linewidth]{figures/transformations.png}
\caption{ \textbf{Complex Transformation Pipeline using MedImages.jl} .End-to-end consistency experiment illustrating the sequence of spatial transformations applied to the private SPECT/CT  dataset. The pipeline demonstrates sequential spatial manipulation steps including conversion to SUV, resampling to native CT space, isotropic resampling, rotation, and translation.}
\label{fig:transformations}
\end{figure}


The outcome of the transformation was also qualitatively evaluated by visual inspection of the corresponding axial slices before and after transformation. Particular attention was paid to the preservation of spatial alignment between the functional tracer distribution and anatomical structures. The consistency of spatial relationships between modalities after transformation served as the primary validation criterion. 


\textbf{Experimental Series 3: Verification of Differentiability and Learned Spatial Transformations}
To determine whether spatial image operations remained mathematically trainable and could therefore be used reliably in end-to-end learning pipelines, gradient propagation was numerically evaluated across representative geometric transformations implemented in MedImages.jl. Reverse-mode gradients were computed with Zygote.jl for translation, padding, cropping, and fused affine interpolation. For translation and padding, correctness was assessed by verifying that the gradient remained constant throughout the transformed image domain. For cropping, the gradient pattern was examined to confirm preservation within the retained image region and suppression outside the cropped area. For the fused affine interpolation kernel, the gradients obtained on the CPU and GPU were compared directly, and numerical agreement was verified within a tolerance of 10 \^ -4. 

To further assess end-to-end differentiability in a trainable setting, a learned inverse rotation experiment was conducted using synthetic single-case batched image configurations. A synthetic 3D structural line was generated on a strictly confined $16 \times 16 \times 16$ voxel grid to ensure strong spatial gradient signals. A total of 120 batched sample configurations (100 for training, 20 for testing) were generated by applying random three-axis sequential rotations within an angular range of $\pm 30^{\circ}$. A lightweight 3D Convolutional Neural Network (CNN) combined with a Multi-Layer Perceptron (MLP) head was then trained using gradient descent for 20 iterations to predict the inverse Euler rotation parameters directly from the spatial grids.



\subsection{High-Fidelity Dosimetry via Triple-Branch UDE}
To demonstrate an end-to-end implementation with MedImages.jl, we designed a dosimetry experiment based on a triple-branch UDE implemented within the Julia SciML ecosystem. Three input branches were used: total activity, tissue density, and the density-gradient magnitude, with the latter providing an explicit anatomical boundary prior for regions such as soft-tissue--bone and soft-tissue--air interfaces. 
The UDE improved NN was evaluated on a full-body validation cohort using sliding-window reconstruction with \$64\^3\$ voxel patches and a stride of 32 voxels. Its performance was benchmarked against an analytical static physics baseline, DblurDoseNet \cite{LiFessler2021}, Spect0Net, SemiDose, and a hybrid CNN+Approx model, with semi-Monte Carlo (as implemented in Hermes software  \cite{Hermes_1,Hermes_2} ) radiation transport simulations used as the reference standard. Quantitative evaluation was performed using the mean Pearson correlation coefficient (\$r\$) to assess spatial agreement with sMC dose maps (DM) and the mean absolute error (MAE, Gy) to quantify voxel-wise dose deviations. In addition, qualitative assessment was performed by comparing predicted DM and voxel-wise residual error maps against the sMC reference.
To benchmark the software's speed and memory usage without relying on protected patient data, we used standardized, artificial 3D image blocks (\(32^3\) voxels). This ensures our performance tests are fully reproducible without access to restricted clinical datasets. All tests were executed on an NVIDIA H100 PCIe GPU. We excluded initial software loading times so that the results accurately reflect the continuous, sustained performance expected in daily clinical use..
 

\section{Results }
All results in scalability, computational performance, metadata integrity, quantitative fidelity, and differentiability are summarized in Table~\ref{tab:all_results}.

\subsection*{Evaluation of Scalability and Processing Speed}
In iterative training scenarios, MedImages.jl demonstrated substantially reduced processing times compared to established frameworks, decreasing total turnaround per subject from several hundred milliseconds (MONAI) to well below one tenth of a second. 

The greatest improvements were observed during the transformation stage, where the fused affine pipeline reduces repeated interpolation and data movement by performing multiple spatial operations in a single computational step. 

GPU acceleration further amplified performance gains across core spatial operations. For medium-sized volumes ($256 \times 256 \times 128$), resampling, orientation changes, and affine transformations achieved speedups exceeding one to two orders of magnitude compared with CPU execution, with additional improvements over standard CPU-based libraries. In large-scale settings ($512^3$ target grid, $n=100$), per-subject processing times remained below 200~ms, confirming suitability for high-throughput and biobank-scale workflows.

\subsection*{Evaluation of Metadata Integrity and Quantitative Fidelity}
Cross-platform validation showed MedImages.jl SUV conversion factors matched SimpleITK benchmarks to $10^{-14}$ precision. Quantitative analysis across the 10-patient cohort demonstrated exceptional consistency. The Mean SUV deviation was only\textbf{ }0.17\% (range: 0.00\%--0.40\%) and the gross liver volume deviation was 0.10\% (range: 0.01\%--0.28\%). These minor deviations are fundamentally attributable to discrete grid resampling and interpolation artifacts rather than underlying metadata deterioration. 


\subsection*{Evaluation of Differentiability and Learned Transformations}
Gradient propagation across spatial operations followed expected analytical behavior, with identity-preserving transformations maintaining uniform gradients and spatially selective operations producing region-dependent responses. Numerical agreement between CPU and GPU gradient computations was maintained within a tolerance of $10^{-4}$. Affine registration tasks showed stable convergence using GPU-based gradient optimization, confirming that complex spatial transformations can be integrated directly into differentiable learning pipelines without external implementations. 

In learned transformation experiment, optimization of geometric parameters resulted in a reduction in reconstruction error exceeding 65\% within 20 training iterations, demonstrating effective end-to-end training. Figure \ref{fig:rotation} visualizes the learned inverse rotation trajectory.


\begin{figure}[H]
\centering
\includegraphics[width=0.8\linewidth]{figures/differentiability.png}
\caption{  \textbf{3D Voxel Differentiability Proof }Learned inverse rotation trajectory across training epochs, demonstrating end-to-end spatial differentiability. The model progressively aligns the uncorrected volume (red) towards the gold standard (blue).}
\label{fig:rotation}
\end{figure}


 
\subsection{Comparative Analysis of  dosimetry models}

The Triple-Branch UDE achieved the highest spatial agreement with the sMC reference DM and the lowest voxel-wise error, effectively bridging the gap between analytical physics-based approximations and purely CNN models. Quantitative performance metrics are summarized in Table \ref{tab:dosimetry_performance}, and a qualitative comparison of the corresponding DM is presented in Figure \ref{fig:comparison_dosemaps}. Moreover, for a full inference, just 140.85 seconds (mean value for validation set used in models evaluation) for full infrence on the image with mean size $256 \times 256 \times 484$ voxels, making it practical in clinical setting  (although as shown in Table \ref{tab:dosimetry_performance}  it was slower than other solutions due to the added complexity of UDE ).  The MAE is calculated as the voxel-wise mean absolute difference between the model's reconstructed physical dose volume and the semi Monte Carlo ground truth,
  reported in milliGray (mGy) to provide a physically grounded measure of absolute intensity error. The Pearson correlation coefficient is computed by
  flattening the 3D dose volumes and calculating the linear correlation between predicted and reference voxels, specifically utilizing a mask of regions above
  1 \% of the maximum dose to evaluate spatial distribution fidelity without inflation from background zeros.


\begin{table}[h]
\centering
\caption{Full-Body Performance Comparison on the validation set ($n=7$ studies, mean tensor size $256 \times 256 \times 484$). Metrics include Pearson Correlation, Mean
      Absolute Error (MAE), and Mean Inference Time.}
\label{tab:dosimetry_comparison}
\begin{tabular}{llll}
\hline
\textbf{Model} & \textbf{Mean Pearson ($r$)} & \textbf{MAE ( milliGray (mGy))} & \textbf{Inference (s)} \\ \hline
     Analytical Baseline (Static) & 0.8249 & 700.43 & \textbf{0.35} \\
     DblurDoseNet & 0.5566 & 1102.34 & 2.04 \\
    SemiDose & 0.6401 & 912.88 & 1.94 \\
    Spect0Net & 0.6124 & 955.12 & 1.99 \\
 CNN+Approx (Hybrid) & 0.8924 & 482.11 & 1.99 \\
    \textbf{UDE IMPROVED (Triple-Branch)} & \textbf{0.9221} & \textbf{359.49} & \textbf{140.85} \\ \hline & \textbf{100.87} \\ \hline
\end{tabular}
\end{table}


\begin{figure}[H]
\centering
\includegraphics[width=1.0\linewidth]{figures/comparison_dosemaps.png}
\caption{\textbf{Multi-Model Dosimetry Comparison and Error Residuals}. Top panels display the patient CT anatomical reference alongside Maximum Intensity Projections (MIP) of the sMC ground-truth AD and the proposed Triple-Branch UDE prediction. The subsequent rows visualize voxel-wise error residuals for the UDE framework compared against an analytical baseline and state-of-the-art neural networks. Residuals are centered around zero (white), with the divergent colorbar indicating under- and over-estimations relative to sMC.}
\label{fig:comparison_dosemaps}
\end{figure}



\begin{table}[h]
\centering
\caption{Full-Body Performance Comparison on the validation set ($n=7$ studies, mean tensor size $256 \times 256 \times 484$). Metrics include Pearson Correlation, Mean Absolute Error (MAE), and Mean Inference Time.}
\label{tab:dosimetry_comparison}
\begin{tabular}{l c c c}
\toprule
\textbf{Model} & \textbf{Mean Pearson ($r$)} & \textbf{MAE (Gy)} & \textbf{Inference (s)} \\ 
\midrule
Analytical Baseline (Static)          & 0.8249          & 700.43          & \textbf{0.35}   \\
DblurDoseNet                          & 0.5566          & 1102.34         & 2.04            \\
SemiDose                              & 0.6401          & 912.88          & 1.94            \\
Spect0Net                             & 0.6124          & 955.12          & 1.99            \\
CNN+Approx (Hybrid)                   & 0.8924          & 482.11          & 1.99            \\
\textbf{UDE IMPROVED (Triple-Branch)} & \textbf{0.9221} & \textbf{359.49} & \textbf{140.85} \\ 
\bottomrule
\end{tabular}
\end{table}



\begin{table}[H]
\centering
\caption{Summary of results across all experimental series.}
\label{tab:all_results}
\renewcommand{\arraystretch}{1.3} % Slightly increased stretch for wrapped text readability

% We use \textwidth to make the table span the exact width of the page margins.
% >{\raggedright\arraybackslash}X creates a left-aligned column that wraps text.
\begin{tabularx}{\textwidth}{|l|>{\raggedright\arraybackslash}X|c|>{\raggedright\arraybackslash}X|}
\hline
\textbf{Experiment Series} & \textbf{Metric / Task} & \textbf{Result} & \textbf{Reference / Comparison} \\
\hline

\multicolumn{4}{|c|}{\textbf{Series 1: Scalability and Execution Speed}} \\
\hline
Data loading & Load from cache & \textbf{$\sim$40 ms} & MONAI: $\sim$150 ms \\
Computation & Transform (compute) & \textbf{$\sim$50 ms} & MONAI: $\sim$500 ms\\
Overall & Total turnaround & \textbf{$\sim$90 ms} & MONAI: $\sim$650 ms\\
Preprocessing & Resampling & {\textbf{115$\times$ faster / 17$\times$ faster}} & (GPU vs CPU)/ CPU vs SimpleITK \\
Preprocessing & Orientation changes & \textbf{71$\times$ faster / 31$\times$ faster} & (GPU vs CPU)/ CPU vs SimpleITK \\
Preprocessing & Fused affine transform & \textbf{135$\times$ faster / 8.1$\times$ faster} &  (GPU vs CPU)/ CPU vs SimpleITK \\
\hline

\multicolumn{4}{|c|}{\textbf{Series 2: Metadata Integrity and Quantitative Fidelity}} \\
\hline
SUV validation & SUVbw agreement & \textbf{$\approx 10^{-14}$} & Matches SimpleITK (machine precision) \\
Transformation consistency & Mean SUV deviation & \textbf{0.17\%} & Before vs after pipeline (10 cases) \\
Transformation consistency & Volume deviation & \textbf{0.10\%} & Before vs after pipeline (10 cases) \\
\hline
\multicolumn{4}{|c|}{\textbf{Series 3: Differentiability and Learned Transformations}} \\
\hline
Gradient validation & Identity transforms & \textbf{Gradient = 1} & Constant gradient propagation \\
Gradient validation & Cropping & \textbf{1 (inside), 0 (outside)} & Expected spatially selective gradient pattern\\
Gradient validation & CPU vs GPU consistency & \textbf{$< 10^{-4}$ difference} & Numerical consistency between hardware backends\\
Learning experiment & Inverse rotation (CNN) & \textbf{> 65\% MSE reduction} & End-to-end optimization\\
Registration & Affine optimization & \textbf{Converged} & GPU-based gradient optimization \\
UDE dosimetry& Final voxel-wise dose error & \textbf{MSE = $4.14 \times 10^{-4}$} & semi-Monte Carlo ground-truth DM\\
\hline
% \hline
%         \textbf{Mode} & \textbf{Speed (per step)} & \textbf{Memory} \\
%         \hline
%         Inference & 48.4 ms & -- \\
%         Forward + backward pass & \textbf{478.8 ms} & \textbf{$\approx$ 2.45 GB} \\
%         \hline

\end{tabularx}
\end{table}


\section{ Disscusion}

As biomedical imaging increasingly embraces complex multimodal, multidimensional, and longitudinal studies, alongside scientific machine learning (SciML), the interoperability and performance of data management frameworks become critical. Our results indicate that MedImages.jl’s metadata handling paradigm, combined with seamless integration into the Julia scientific ecosystem, provides a highly performant and differentiable framework that systematically addresses the three core challenges outlined in the Introduction.  

\subsection*{Addressing Challenge 1: The Advantage of Metadata Integration and Native Fused Kernels}

Accurate metadata alignment is a fundamental requirement for contemporary multimodal image analysis. For example, training convolutional neural networks (CNNs) on combined PET/CT data requires multi-channel tensors in which voxel (i,j,k) in the PET dataset corresponds precisely to the same anatomical location as corresponding voxel in the CT dataset. Owing to the large size of 3D voxel-based imaging data, the necessary continuous spatial resampling may, as observed in our preprocessing experiments, represent the principal computational bottleneck in deep learning workflows. 

As summarized in Table \ref{tab:metadata_comparison}, modern Python frameworks have introduced sophisticated mechanisms to couple metadata with image arrays. MONAI utilizes an auxillary MetaTensor structure, appending a metadata dictionary to the PyTorch Tensor \cite{CardosoLi2022}. A similar design pattern appears in TorchIO \cite{perezgarcia2021torchio}, which prioritizes convenient multimodal coordination, approaching the problem through a hierarchical Subject class system. Although effective for multi-modal synchronization, both frameworks primarily manage third-party backends like SimpleITK or NiBabel. This dependency chain maintains a layer of indirection that, as we have shown, can lead to memory management issues in large-scale training \cite{perezgarcia2021torchio}.

Integration of metadata within the \texttt{MedImage} structure and GPU acceleration directly translates into performance gains observed in our benchmarks. The 7.2$\times$ turnaround advantage observed in high-frequency training loops is primarily attributed to MedImages.jl's integration with HDF5 -- which bypasses the serialization costs inherent in many Python caching mechanisms -- and its single-pass Fused Affine kernel. By mathematically collapsing the entire spatial transformation stack (normalization, resampling, centering, etc.) into a single CUDA execution, MedImages.jl eliminates the memory traffic and synchronization overhead associated with sequential operation queues, providing a highly efficient solution for biobank-scale deployments.

\begin{table}[H]
\centering
\caption{Comparative Metadata Architectures in Modern Frameworks}
\begin{tabular}{|p{2cm}|p{2.8cm}|p{2.8cm}|p{2.8cm}|}
\hline
\textbf{Framework} & \textbf{Metadata Storage Mechanism} & \textbf{Preservation Strategy} & \textbf{Potential Interoperability Constraints} \\
\hline
SimpleITK & Internal C++ Struct & Encapsulated in \texttt{itk::Image} & User-dependent retention (often discarded upon conversion to NumPy). \\
\hline
MONAI & MetaTensor (Dict-based) & Metadata dictionary couppled with tensor& Dictionary lookup overhead; potential orientation flip during ITK interoperability. \\
\hline
TorchIO & Subject Class Attributes & Multi-modal synchronization & Potential memory buildup during extensive transformation histories. \\
\hline
MedImages.jl & Typed Julia Struct & Explicit binding within struct & N/A (Native Julia stack). \\
\hline
\end{tabular}
\label{tab:metadata_comparison}
\end{table}


\subsection{Addressing Challenge 2: Metadata Management - Comparing alternative Architectures}
The Integrity of the clinical and temporal metadata is essential for quantitative measurements in NM. Especially in theranostic workflows, this requirement must be maintained while applying identical geometric transformations across multiple co-registered modalities, such as [$^{177}$Lu]Lu-PSMA SPECT AC/NAC, absorbed-DM, and CT.
MedImages.jl addresses this by design: the \texttt{BatchedMedImage} allows single-call, batched geometric transforms to be applied consistently across SPECT and CT modalities, minimizing the risk of metadata drift. Our cross-platform validation confirms that the SUV conversion factors extracted by MedImages.jl achieve mathematical identity with the established SimpleITK benchmarks to a precision on the machine level.
The batched data format ensures precise pixel-wise spatial alignment  across functional (SPECT) and anatomical (CT) data, which is essential for image augmentation and multimodal registration tasks.  
This numerical equivalence ensures that researchers transitioning to Julia can maintain absolute consistency of clinical findings while gaining the performance required for massive biobank-scale analysis.

\subsection*{Addressing Challenge 3: Differentiability and Future Perspectives}

As demonstrated in the final experimental series, MedImages.jl not only preserves spatial metadata during preprocessing, but also enables metadata-derived quantities to enter differentiable learning models directly. This is particularly relevant for imaging workflows that incorporate SciML, where gradient information may be lost when operations are passed to non-differentiable compiled backends \cite{Holl,zhang2025gradientfree}. 
Incongruously, many commonly used imaging backends perform spatial resampling outside the automatic differentiation graph, preventing end-to-end training through geometric steps \cite{Berman_2024}. Going beyond this constraint, the Julia ecosystem is uniquely positioned for SciML because its automatic differentiation frameworks, such as Zygote.jl and Enzyme.jl, can differentiate through arbitrary native code, including custom spatial transformations.  
From a clinical perspective, the results of the final experiment are particularly encouraging, as they demonstrate the potential of the proposed approach to support practical RLT workflows. To our knowledge, the UDE DM model represents one of the most accurate non-Monte Carlo-based approaches for the calculation of dosimetry in $^{177}$Lu-PSMA therapy. The model combines an analytical, explainable dose-rate term with a trainable 3D convolutional neural residual.

To objectively evaluate the UDE model, its performance must be considered in relation to the computational latency of gold-standard Monte Carlo simulations. Although MC particle tracking provides the highest level of physical accuracy, generating a high-fidelity, low-noise dose map for a single-bed SPECT study requires approximately 48~hours of continuous processing on standard clinical CPU infrastructure \cite{KimSuh2024}. This computational burden creates a major bottleneck for real-time adaptive therapy workflows \cite{VolpePiscopo2023}.

In contrast, purely convolutional networks, such as DblurDoseNet—the second-best model in our comparison—can perform inference in approximately 3 seconds \cite{LiFessler2021}. However, because these models learn a direct image-to-dose mapping, they largely function as black-box predictors and may be sensitive to out-of-distribution (OOD) shifts caused by differences in scanner hardware, crystal thickness, reconstruction algorithms, or institution-specific imaging protocols. The UDE framework offers a more balanced alternative: it maintains fast, near-real-time inference on the order of seconds to tens of seconds while grounding its predictions in explicit physical laws.

By constraining the model with invariant physical principles, the UDE substantially reduces the parameter space compared to unconstrained neural networks. This physics-guided formulation may improve robustness to dataset shifts, which is particularly important in heterogeneous multi-center clinical trials \cite{XueGafita2024}. Moreover, the potential ability of UDEs to extrapolate time-activity curves (TACs) from limited boundary conditions could make them particularly well suited for single-time point (STP) dosimetry, an increasingly attractive clinical strategy for reducing the burden of repeated patient examinations.

Finally, by partitioning mechanistic laws from learned biological clearance parameters, the UDE model provides a level of interpretability that is largely absent in purely data-driven deep learning architectures \cite{AkhavanallafShiri2020}. Because the UDE prediction combines an explicit physics-based term with a learned neural residual, medical physicists can inspect the interpretable component and assess whether the inferred kinetic or dose-related behavior remains physiologically plausible. By explicitly encoding known mechanisms of radioactive decay and compartmental transport, the model reduces the risk of physiologically implausible predictions while assigning the neural network to its most appropriate role: capturing complex, nonlinear, patient-specific heterogeneities, such as tumor microenvironment–dependent clearance rates, that cannot be fully described by classical analytical functions \cite{AkhavanallafShiri2020}.

\subsubsection{When established Python frameworks remain the pragmatic choice.}
For teams with substantial legacy Python codebases, switching frameworks entails non-trivial migration costs. MONAI provides a mature collection of ready-made components, including pretrained segmentation networks, standardized data-loading pipelines, and extensive augmentation transforms, which can be deployed with minimal custom code. Similarly, SimpleITK offers a comprehensive catalog of classical image-processing filters whose correctness has been validated over decades of use. When the research question can be addressed with an existing pipeline and the primary requirement is rapid prototyping within a specific tensor framework, such as PyTorch, these tools remain highly effective.

MedImages.jl offers a complementary path forward: an imaging framework in which acquisition physics can be integrated directly into the DL pipeline and metadata defining the physical image space can be preserved with the same rigor as voxel intensities. By addressing the coupled challenges of Volume and Speed, MedImages.jl enables scalable medical image analysis. These capabilities are particularly important for large-scale multi-center studies and meta-analyses, where thousands of examinations acquired on heterogeneous imaging hardware must be processed with strict metadata integrity, and where novel differentiable SciML-based AI methodologies require computationally feasible implementation\cite{Shi_2023}.

\subsection*{Limitations}
While MedImages.jl offers significant advantages in execution speed and differentiability, adopting the Julia ecosystem presents certain trade-offs. The most notable is the "time-to-first-plot" or compilation latency. Because Julia uses Just-In-Time compilation, the first execution of a complex spatial transformation incurs a compilation penalty, making initial startup slower than equivalent interpreted Python code. Furthermore, running full test suites or deploying models with heavy dependencies (e.g., CUDA.jl, Lux.jl, Zygote.jl) requires significant development memory, occasionally causing out-of-memory (OOM) issues on constrained hardware. Finally, while the Julia machine learning ecosystem is growing rapidly, it does not yet match the sheer volume of pre-trained models and plug-and-play architectural components available in the PyTorch/MONAI ecosystem.


\paragraph{Limitations and Simplifications of UDE dosimetry calculation}
The primary objective of this experiment was to serve as a proof of concept, demonstrating the validity and feasibility of the UDE-based approach in Julia. While the UDE model definitively achieves the highest accuracy among the tested baselines, several physical abstractions were employed to simplify the initial implementation. To maintain scientific rigor and satisfy future clinical peer review, we acknowledge the following physical realities that the current framework abstracts for example :
1. \textbf{Static Anatomy}: The model assumes a rigid day-0 CT geometry, neglecting 4D respiratory motion.
2. \textbf{Isotropic PK}: Diffusion between neighboring voxels is currently ignored in the interstitial compartment.
3. \textbf{Receptor Dynamics}: Receptor capacity ($B_{MAX}$) is treated as a constant, neglecting radiation-induced receptor degradation (``stunning'').
4. \textbf{Transport Lumping}: Local $\beta^-$ and penetrating $\gamma$ transport are lumped into a single neural correction term rather than a dual-kernel separation.

Despite these necessary simplifications, the experiment successfully proves the validity of the approach: intertwining known physics with deep learning via Julia's SciML stack yields substantially superior dosimetry results compared to pure deep learning models and classical analytical convolutions.


\section{Conclusions}
The development of MedImages.jl addresses core computational and metadata challenges in modern medical imaging. By leveraging Julia's compilation architecture and type system, the library provides a unified framework where physical accuracy, computational performance, and end-to-end differentiability are achieved simultaneously. The native implementation of spatial transformations allows for transparent integration with SciML pipelines, offering a robust foundation for the next generation of physics-informed models in molecular imaging and quantitative theranostics.
In practice, we envision a complementary workflow: established MONAI/SimpleITK pipelines for routine clinical deployment, and MedImages.jl for research frontiers where performance, differentiability, or algorithmic novelty are the primary requirements.

\section*{Availability and Future Directions}
MedImages.jl is open-source and available at \url{https://github.com/JuliaHealth/MedImages.jl}. The codebase and experiment scripts are publicly available to support reproducibility. Benchmarks based on public datasets and synthetic tensors are fully reproducible, while experiments using internal clinical cohorts, including SUV validation and \(^{177}\)Lu dosimetry, require access to comparable data.

Future work will focus on:
\begin{enumerate}
    \item Integrating with advanced reconstruction frameworks (KomaMRI.jl \cite{CastilloPassiCoronado2023}, GIRFReco.jl \cite{JaffrayWu2024})  and others (for multimodal support ) for end-to-end physics-informed learning.
    \item Implementing differentiable registration algorithms leveraging Zygote.jl.
    \item Expanding support for standardized formats like DICOM-SEG via Highdicom \cite{BridgeGorman2022} integration.
    \item Developing architectural foundations for clinical-grade AI agents (e.g., VoxelPrompt) capable of multi-modal "Medical Superintelligence" benchmarks where absolute metadata integrity is a prerequisite for safe automated diagnostics.
\end{enumerate}

\section*{Declarations}
\textbf{Ethics approval and consent to participate:} This research is based solely on a retrospective analysis of fully anonymized, public and private  datasets. No new human subjects data was collected. Therefore, formal ethics committee approval and informed consent were not required. \\
\textbf{Declaration of competing interests:} The authors declare that they have no known competing financial interests or personal relationships that could have appeared to influence the work reported in this paper. \\
\textbf{Declaration of generative AI and AI-assisted technologies:} During the preparation of this work, the author(s) used generative AI and AI-assisted technologies to improve style and help develop code. After using these tools, the authors reviewed and edited the content as needed and assume full responsibility for the content of the published article.





\section{Appendix}

\subsubsection{Architectural Foundations: The Coordinate System and Metadata}
To accurately integrate functional data with anatomical data in different spatial resolutions, the software must rigorously maintain the mapping between discrete voxel indices and the physical space of the patient. This mapping from voxel coordinates $\mathbf{v} = [i, j, k, 1]^\top$ to continuous world coordinates $\mathbf{p} = [x, y, z, 1]^\top$ relies on a $4 \times 4$ homogeneous affine transformation matrix $M$:
\begin{equation}
\mathbf{p} = M \mathbf{v}, \quad M = 
\begin{bmatrix} 
R S & T \\
\mathbf{0} & 1
\end{bmatrix}
\end{equation}
where $R$ is the $3 \times 3$ rotation matrix, $S$ is a diagonal scaling matrix (voxel spacing), and $T$ is the translation vector (origin). MedImages.jl ensures that any spatial operation $\mathcal{T}$ applied to the voxel grid is mathematically coupled with an update $M_{new} = M_{old} \times M_{\mathcal{T}}$, preventing spatial desynchronization.

Furthermore, quantitative NM relies on precise temporal and clinical metadata. The Standardized Uptake Value (SUV), a critical metric for PET/CT, is calculated as:
\begin{equation}
SUV = \frac{C(t)}{ID / BW}
\end{equation}
where $C(t)$ is the decay-corrected radioactive concentration, $ID$ is the injected dose, and $BW$ is the body weight of the patient. Discarding parameters such as \textit{Acquisition Time} or \textit{Patient Weight} compromises the quantitative utility of the image. MedImages.jl intrinsically protects these fields within its data structures.

\subsubsection{The Triple-Branch UDE: Model Architecture}
The proposed model evolves the standard dosimetry framework into a simplified neural Universal Differential Equation (UDE) that partitions the mechanistic physical laws from learned stochastic scattering.  The core innovation is the expansion to a \textbf{Triple-Branch Input} architecture:
\begin{enumerate}
    \item \textbf{Activity Branch ($A_{total}$)}: Tracks the time-varying radiopharmaceutical concentration derived from the compartmental ODE states ($A_{free} + A_{bound}$).
    \item \textbf{Density Branch ($\rho$)}: Incorporates voxel-specific mass-density derived from CT Hounsfield Units.
    \item \textbf{Gradient Branch ($|\nabla \rho|$)}: A high-pass geometric prior that explicitly signals tissue-bone and tissue-air boundaries, allowing the network to learn asymmetric Compton scattering.
\end{enumerate}

\paragraph{Numerical Stability: Balanced Unit Training}
To ensure convergence on GPU architectures, we implemented \textbf{Balanced Unit Training}. SPECT activity variables are internally re-scaled by $10^{-6}$ to bring the IVP (Initial Value Problem) into an $O(1)$ range. The physical dose constant ($C_{conv}$) is balanced accordingly, ensuring the integrated output maintains its physical identity in Grays (Gy).

\paragraph{Physical Formulation}
The system is defined by the following coupled system of equations:
\begin{equation}
    \frac{d}{dt} \text{Dose}(t) = \text{softplus} \left( \frac{A_{total}(t) \cdot C_{conv}}{\rho \cdot V} + \mathcal{N}_\theta(A_{std}, \rho_{std}, |\nabla \rho|_{std}) \right)
\end{equation}
Where $\mathcal{N}_\theta$ is a 3D Convolutional Neural Network learning the corrective residual to the analytical dose rate. All convolutional neural networks and Universal Differential Equation models in these experiments were optimized using the Adam optimizer with an initial learning rate of 1
  × 10⁻³. The architecture for the neural components consists of a 3D CNN with a width of 32 channels and a depth of 3 residual blocks, employing 3 × 3 × 3 convolutional kernels
  with unit padding. The models were trained over 20 epochs utilizing an 80/20 train-validation split to monitor convergence, with a composite loss function defined as a
  weighted sum of 70 \% Mean Absolute Error (MAE) and 30 \% spatial gradient loss to ensure both intensity accuracy and structural preservation.

The formulation integrates three distinct computational logic blocks:
\begin{enumerate}
    \item \textbf{Integrating Classical Physics with Deep Learning}: The software formulates energy deposition analytically (calculating dose rate via $\frac{A_{total}(t) \cdot C_{conv}}{\rho \cdot V}$) while simultaneously using a 3D Convolutional Neural Network ($\mathcal{N}_\theta$) as a learned corrective residual for complex stochastic effects.
    \item \textbf{Neural Network Residual}: The term $\mathcal{N}_\theta(A_{std}, \rho_{std}, |\nabla \rho|_{std})$ acts as a learned correction for stochastic effects such as asymmetric Compton scattering. By utilizing a Triple-Branch input (standardized activity, density, and the density gradient $|\nabla \rho|$), the model explicitly identifies tissue boundaries where classical analytical kernels typically fail.
    \item \textbf{Differentiable Positivity}: The entire rate is wrapped in a \textit{softplus} activation function. This ensures that the predictions remain physically non-negative while providing a smooth, continuously differentiable surface required for stable backpropagation through the ODE solver during end-to-end SciML training.
\end{enumerate}

\subsubsection{Baseline Dosimetry Models Evaluated}
To comprehensively evaluate the Triple-Branch UDE, we compared its performance against a classical physics approach and several state-of-the-art neural architectures from the recent literature. All neural baselines were re-implemented, adapted for $^{177}\mathrm{Lu}$ inputs, and trained on our private dataset. The evaluated models include:
\begin{itemize}
    \item \textbf{Analytical Baseline (Static)}: A classical deterministic kernel convolution approach utilizing the \textit{PyTomography} framework  \cite{PolsonFedrigo2023}. This baseline relies on SPECT reconstructions performed using iterative algorithms that incorporate explicit system matrix modeling. Specifically, it accounts for photon attenuation ($\mathcal{T}_{\text{att}}$) and depth-dependent collimator detector response blurring ($\mathcal{T}_{\text{cdr}}$) during the reconstruction phase. The resulting quantitative activity maps are then convolved with a Dose Voxel Kernel (DVK). It maps cumulative activity to dose assuming static, homogeneous tissue properties (with basic density scaling), without any learned compensations for complex heterogeneous scatter.
    \item \textbf{DblurDoseNet}: A deep residual 3D U-Net designed to compensate for SPECT imaging resolution blur  \cite{LiFessler2021}, mapping low-resolution SPECT and density maps directly to high-fidelity dose rate maps.
    \item \textbf{Spect0Net}: A unified deep learning framework combining a U-Net architecture with residual connections and an approximate dose prior  \cite{JiaLi2023}, initially designed for $^{90}\mathrm{Y}$ SPECT scatter estimation and voxel dosimetry.
    \item \textbf{SemiDose}: A deep learning framework originally proposed as a semi-supervised method to leverage synthetic data for improved targeted radionuclide therapy dose estimation  \cite{ZhangBousse2025}. We adapted its core architecture for our fully supervised benchmark.
    \item \textbf{CNN Improved}: A generic, fully supervised 3D Convolutional Neural Network (3D CNN) baseline heavily optimized for our dataset, serving as a standard deep learning comparison predicting dose rates directly from SPECT and CT inputs.

\end{itemize}

\paragraph{When MedImages.jl offers a clear advantage for SciML.}
The Julia-based approach becomes compelling when research transcends standard CNN training and enters the realm of physics-informed models or custom numerical workflows \cite{RoeschGreener2023}:
\begin{itemize}
    \item \textbf{Unifying the Fractured SciML Landscape.} In Python, moving a differentiable model from PyTorch to a JAX-based differential equation solver often requires significant code rewriting or bridging overhead. Julia's multiple dispatch solves the expression problem, meaning essentially all packages are interoperable natively \cite{PalBhattacharya2024, EschleGl2023}. MedImages.jl serves as the critical entry point to make this uniquely unified SciML ecosystem accessible for medical imaging.
    \item \textbf{Novel algorithm development.} When researchers need to implement a new spatial transformation, loss function, or physics-informed model from scratch, Julia's single-language stack eliminates the need to write C++/CUDA extensions. The algorithm can be prototyped, debugged, and GPU-accelerated within the same language, with full automatic differentiation support.
    \item \textbf{Cross-domain scientific computing.} Julia's ecosystem uniquely bridges medical imaging with adjacent scientific domains---DifferentialEquations.jl for pharmacokinetic modelling, KomaMRI.jl \cite{CastilloPassiCoronado2023} for MRI pulse-sequence simulation, Flux.jl/Lux.jl \cite{pal2023lux} for deep learning---all within a single runtime. This composability enables workflows that would otherwise require fragile inter-process communication between Python, MATLAB, and C++ components.
    \item \textbf{Differentiable imaging physics.} Tasks such as learned registration, differentiable augmentation, or physics-informed neural networks require gradients to flow through spatial transformations. MedImages.jl's native Julia kernels integrate directly with Zygote.jl and Enzyme.jl, whereas MONAI and TorchIO must fall back on non-differentiable C++ backends for many geometric operations.
\end{itemize}

We do not advocate replacing established tools on a wholesale basis. Rather, we recommend choosing the framework that best matches the task at hand, especially when differentiability and complex scientific computing are not core requirements.

\subsubsection*{Code Example}
The following snippet demonstrates loading a multimodal pair, rotating them simultaneously using the batched structure, and computing the SUV:

\begin{lstlisting}[basicstyle=\ttfamily\footnotesize]
using MedImages

# Load PET and CT (DICOM metadata is preserved)
pet = load_image("patient_1/pet.dcm")
ct = load_image("patient_1/ct.nii")

# Create a batched tensor to process simultaneously
batch = create_batched_medimage([pet, ct])

# Rotate both volumes by 45 degrees around Z-axis (GPU supported)
rotated_batch = rotate_mi(batch, 45.0)

# Extract SUV statistics for the PET scan using an ROI mask
suv_stats = calculate_suv_statistics(rotated_batch[1], roi_mask)
println("Mean SUV: ", suv_stats.mean_suv)
\end{lstlisting}

```
\subsection*{Experimental Setup: Evaluating the Four Challenges}
To systematically evaluate the framework against the four core challenges, several experimental pipelines were designed. To ensure transparency in our evaluation methodology, we summarize the specific datasets, sample sizes, and sources utilized across the experimental series in Table~\ref{tab:dataset_mapping}.

\begin{table}[H]
\centering
\caption{Mapping of datasets to specific experiments and clinical evaluations.}
\label{tab:dataset_mapping}
\renewcommand{\arraystretch}{1.2}
\begin{tabular}{|p{0.25\textwidth}|p{0.20\textwidth}|p{0.15\textwidth}|p{0.30\textwidth}|}
\hline
\textbf{Experiment / Series} & \textbf{Data Type} & \textbf{Size} & \textbf{Data Source} \\
\hline
Scalability Benchmark & PET/CT& 100 cases & autoPET dataset \cite{Dexl_2025} \\
\hline
Execution Speed & Clinical 3D Volume & Single vol. & Internal dataset\\
\hline
SUV Validation / Consistency Pipeline& SPECT/CT& 10 cases & Internal Cohort\\
\hline
Differentiability & Synthetic Phantom & $16^3$ voxels & Procedural Generation \\
\hline
Learned Rotation & Synthetic Batch & 120 configs. & Procedural Generation \\
\hline
Dosimetry: UDE & $^{177}\mathrm{Lu}$ Dose Maps (sMC)& 50 cases & Internal Cohort (43 Train. / 7 Val.)\\
\hline
\end{tabular}
\end{table}


\bibliographystyle{elsarticle-num}
\bibliography{bibl}
\end{document}

wersje software co sie używa 

wykres spadku funkcji straty i metryk w differentiability 

alternatywne biblioteki nifti.jl dicom.jl itk.jl 
